\documentclass[11pt,a4paper]{article}

\usepackage[utf8]{inputenc}
\usepackage[T1]{fontenc}
\usepackage[ngerman]{babel}
\usepackage{geometry}
\geometry{margin=2.5cm}
\usepackage{amsmath,amssymb,mathtools}
\usepackage{braket}
\usepackage{enumitem}
\setlist[enumerate,1]{label=\textbf{\arabic*.},leftmargin=*,itemsep=.7em}
\setlist[enumerate,2]{label=(\alph*),leftmargin=1.4em,itemsep=.3em}

\title{Klausur — Quantenmechanik}
\author{}
\date{}

\begin{document}
\maketitle

\noindent\textbf{Name:} \rule{0.6\linewidth}{0.4pt}\hfill\textbf{Matr.-Nr.:} \rule{0.25\linewidth}{0.4pt}

\vspace{1em}

\begin{enumerate}

%---------------------- 1 ----------------------
\item \textbf{Geladenes Teilchen im Magnetfeld}\\
Ein Teilchen der Masse \(m\) und Ladung \(q\) bewege sich in einem magnetischen Vektorpotential \(\mathbf A(\mathbf r)\).
\begin{enumerate}
  \item Zeigen Sie, dass dieses System (in geeigneten Koordinaten/Eichungen) durch einen harmonischen Oszillator beschrieben werden kann.
  \item Bestimmen Sie die Energieeigenwerte des Systems.
\end{enumerate}

%---------------------- 2 ----------------------
\item \textbf{Komplexes \(\delta\)-Potential}\\
Betrachten Sie das eindimensionale Potential
\[
V(x) = \bigl(\alpha + \mathrm i\,\gamma\bigr)\,\delta(x)\,,
\]
wobei \(\alpha,\gamma\in\mathbb R\).
\begin{enumerate}
  \item Begründen Sie physikalisch, wie der Ansatz der Wellenfunktion für die Streuung an \(x=0\) aussehen muss.
  \item Bestimmen Sie Reflexions- und Transmissionskoeffizient \(R\) und \(T\).
  \item Gilt \(R+T=1\)? Falls nein, warum nicht?
\end{enumerate}

%---------------------- 3 ----------------------
\item \textbf{Zeitentwicklung im unendlichen Potentialkasten}\\
Ein Teilchen befinde sich im unendlichen Potentialkasten mit \(V(x)=0\) für \(0<x<L\) und \(V(x)=\infty\) sonst. Der Anfangszustand sei
\[
\psi(x,0) \;=\; \Bigl[\sin\!\Bigl(\tfrac{2\pi x}{L}\Bigr) + \cos\!\Bigl(\tfrac{3\pi x}{L}\Bigr)\Bigr]\,
\sin\!\Bigl(\tfrac{\pi x}{L}\Bigr).
\]
\begin{enumerate}
  \item Bestimmen Sie die zum Zeitpunkt \(t=0\) sowie zum Zeitpunkt \(t>0\) möglichen messbaren Energien (Eigenwerte) und die zugehörigen Wahrscheinlichkeiten.
  \item Berechnen Sie den Erwartungswert der Energie \(\langle H\rangle(t)\) sowie die Energieunschärfe \(\Delta E(t)\) zum Zeitpunkt \(t\).
  \item Wie müssen die Erwartungswerte \(\langle x\rangle(t)\) und \(\langle p\rangle(t)\) qualitativ aussehen?
\end{enumerate}

%---------------------- 4 ----------------------
\item \textbf{Endlich-dimensionaler Operator}\\
Es sei
\[
H \;=\; \varepsilon\,\mathbb I \;+\; \lambda\,(A + A^\dagger),\qquad
A \;=\; \ket{2}\!\bra{1} + \ket{1}\!\bra{2} + \ket{2}\!\bra{3}.
\]
% Hinweis: So wie definiert ist A nicht unitär (AA^\dagger \neq \mathbb I). Die Teilaufgabe (a) ist in dieser Form widersprüchlich.
\begin{enumerate}
  \item Zeigen Sie, dass \(A\) unitär ist und dass \(A\) mit \(H\) kommutiert.
  \item Bestimmen Sie die Eigenwerte und Eigenvektoren von \(H\).
\end{enumerate}

%---------------------- 5 ----------------------
\item \textbf{Drehimpuls mit Querfeld}\\
Gegeben sei
\[
H \;=\; H_0 + H_1,\qquad
H_0 \;=\; a\,\mathbf L^2 + b\,L_z,\qquad
H_1 \;=\; c\,L_x,
\]
mit Konstanten \(a,b,c\in\mathbb R\) und den üblichen Drehimpulsoperatoren \(\mathbf L^2, L_x, L_z\).
\begin{enumerate}
  \item Bestimmen Sie die Energieeigenwerte und Eigenzustände von \(H_0\).
  \item Berechnen Sie die Störung der Energie zur führenden Ordnung in \(c\).
  \item Definieren Sie einen geeigneten neuen Operator \(J\) und berechnen Sie den exakten Wert der Energiekorrektur.
\end{enumerate}

\end{enumerate}

\vfill
\noindent\textit{Hinweis:} Sofern nicht anders angegeben, dürfen Sie \(\hbar\) explizit mitführen (keine natürlichen Einheiten).

\end{document}
